% !TEX root = ../main.tex

\sysusetup{
  keywords  = {异常检测, PyTOD, FlashANNS, RAPIDS, FAISS-GPU, GPU 常驻执行, 多 GPU 扩展, 近似最近邻检索, 金融风控},
  keywords* = {Outlier Detection, PyTOD, FlashANNS, RAPIDS, FAISS-GPU, GPU-Resident Execution, Multi-GPU Scaling, Approximate Nearest Neighbor Search, Financial Risk Control},
}

\begin{abstract}
在金融风控、反欺诈、反洗钱及交易监测等实际应用中，异常检测系统面临着双重挑战：既要准确识别极少量的高风险样本，又必须在处理海量数据时满足严格的吞吐率和实时性要求。随着金融交易数据量的不断攀升，系统瓶颈已不再局限于单一检测模型的精度，而是越来越多地体现在异常评分计算复杂、候选召回过程开销大、CPU 与 GPU 间数据交换频繁以及多 GPU 并行扩展效率不佳等系统级问题上\citep{aggarwal2017outlier,zhao2022tod}。传统的CPU方案在处理高维大规模数据时，难以在保证检测效果的同时实现近实时响应。另一方面，若仅对个别算法或检索模块进行GPU加速，而缺乏系统级的协同设计，则候选召回、异常评分与并行扩展等环节之间仍会存在效率瓶颈，无法充分发挥硬件潜力\citep{zhao2022tod,xiao2025flashanns}。

为解决上述问题，本文研究并实现了一个面向金融大规模数据异常检测的并行计算系统，该系统深度融合了PyTOD与FlashANNS。系统中，PyTOD负责高效的异常评分，而FlashANNS则提供高吞吐的候选召回能力。我们在统一的数据流、算子流和调度流框架下，将二者有机结合，构建了一条协同的候选召回-异常评分执行链路\citep{zhao2022tod,xiao2025flashanns}。

针对单GPU场景，我们设计了一条GPU主导的端到端执行链路。该链路基于FlashANNS的异步I/O-计算流水线，致力于让查询特征、候选集以及评分中间张量尽可能驻留在GPU内存中，仅通过pinned memory等机制在必要时与主机进行交互。这一设计有效减少了SSD I/O与GPU计算间的串行等待，也降低了CPU与GPU间不必要的数据搬运开销\citep{xiao2025flashanns,rapids_rmm_pinned,gds_overview}。此外，为了改善原始系统中存在的中间结果序列化及主机频繁参与候选组织的问题，我们引入了RAPIDS对特征处理和表格型中间结果的组织路径进行GPU化重构，并利用FAISS-GPU对候选检索后端进行了工程强化\citep{rapids_cudf_pandas,faiss_gpu_wiki}。

在多GPU扩展方面，我们设计并实现了基于数据并行与索引复制的方案。系统优先采用replicated架构，使得每块GPU都能独立完成本地的候选召回与异常评分，而CPU则主要肩负控制面调度与小规模最终结果汇总的职责\citep{cuvs_multi_gpu_nn,cuvs_dist_mode}。当索引规模或业务约束使得跨卡结果归并不可避免时，我们将归并逻辑尽可能下沉到GPU侧执行，并利用NCCL进行设备间通信。同时，结合拓扑感知的I/O通道绑定与准入控制机制，以减轻共享存储路径对多GPU扩展效率的影响\citep{nccl_stream_semantics,gds_design_guide}。

实验部分，我们使用1M规模的现实与半现实数据集验证了方法的有效性，并构建了10M、50M等规模的高保真合成数据集，用于测试系统的吞吐量、延迟及扩展能力\citep{ieee_cis_fraud,ibm_aml_data,dgraphfin_baseline}。实验结果表明，我们所提出的PyTOD--FlashANNS融合系统在保证检测效果可接受的前提下，能显著提升整体吞吐率，并有效降低CPU--GPU间的数据往返与阶段同步开销。在多GPU实验中，数据并行与索引复制策略也为系统扩展提供了良好支撑\citep{xiao2025flashanns,cuvs_multi_gpu_nn}。

需要说明的是，本研究的目标并非提出新的异常检测理论模型，而是立足于金融大规模数据处理这一具体场景，对候选召回、异常评分、数据路径组织及多GPU扩展等问题进行系统化的工程设计与实现。研究结果表明，通过将PyTOD与FlashANNS深度融合，能够有效缓解异常检测流程中计算开销大、存储访问代价高和并行扩展受限等核心矛盾。这为在金融风控等场景中构建高吞吐、近实时的异常检测系统，提供了一条具备工程可行性的技术路径\citep{zhao2022tod,xiao2025flashanns}。
\end{abstract}

\begin{abstract*}
Outlier detection systems in financial risk control applications—including fraud detection, anti-money laundering, and transaction monitoring—face a dual challenge: they must accurately identify a small number of high-risk cases while simultaneously meeting strict throughput and latency requirements at scale. As transaction volumes grow, the primary bottleneck has shifted from the accuracy of individual detection models to system-level issues. These include computationally expensive anomaly scoring, costly candidate retrieval, frequent CPU–GPU data movement, and limited multi-GPU scalability.

To tackle these challenges, we developed a parallel-computing-based \emph{PyTOD--FlashANNS} integrated outlier detection system for large-scale financial data. In this system, PyTOD acts as the primary scoring backbone, while FlashANNS provides high‑throughput candidate retrieval. Within a unified execution pipeline, candidate retrieval and anomaly scoring are coordinated to work in synergy, enabling the system to improve end‑to‑end processing efficiency without compromising detection effectiveness.

For single‑GPU scenarios, we designed a GPU‑dominant end‑to‑end execution path. Leveraging FlashANNS's asynchronous I/O–compute pipeline, query features, retrieved candidates, and intermediate scoring tensors are retained on the GPU as much as possible; host‑side interaction is confined to necessary boundaries via mechanisms like pinned memory. This approach reduces both serialized waiting between SSD I/O and GPU computation as well as unnecessary CPU–GPU data transfers. Furthermore, we introduced RAPIDS to GPU‑accelerate feature processing and tabular intermediate organization, while employing FAISS‑GPU to engineer a more efficient retrieval backend.

For multi‑GPU scaling, we implemented an extension scheme based on data parallelism and index replication. The system adopts a replicated architecture as the default, allowing each GPU to independently perform local candidate retrieval and anomaly scoring, while the CPU primarily handles control‑plane scheduling and final aggregation of small result sets. When index capacity or deployment constraints necessitate cross‑GPU merging, the merge logic is offloaded to the GPU side using NCCL for communication. Topology‑aware I/O lane binding and admission control are also employed to mitigate the impact of shared storage paths on multi‑GPU scalability.

Experiments using 1M‑scale real and semi‑real datasets, along with high‑fidelity synthetic datasets at 10M and 50M scales, demonstrate that the proposed system substantially improves overall throughput while reducing CPU–GPU round‑trip traffic and synchronization overhead. In multi‑GPU tests, the combination of data parallelism and index replication offers a practical foundation for system scaling.
\end{abstract*}
